% Part 3 – Special features: Code snapshots + Generative AI prompts
% เฉพาะ 4 feature ที่เลือก: AI word-guessing, RPS, File/picture, Discord bridge (ไม่มี Cloud กับ /ping)
% สร้างจากโค้ดจริง (เลขบรรทัดตรงกับไฟล์ ณ วันที่ 2026-10-28)
% Compile ด้วย XeLaTeX (ไม่ใช่ pdfLaTeX) เพราะมีภาษาไทย
%   Overleaf: Menu -> Compiler -> XeLaTeX
%   ฟอนต์ Norasi / TlwgMono มากับ TeX Live (แพ็กเกจ fonts-tlwg) ถ้าไม่มีให้เปลี่ยนเป็นฟอนต์ไทยอื่น
\documentclass[11pt,a4paper]{article}

\usepackage[margin=2cm]{geometry}
\usepackage{fontspec}
\setmainfont{Norasi}
\setmonofont{TlwgMono}[Scale=0.95]
\XeTeXlinebreaklocale "th"
\XeTeXlinebreakskip = 0pt plus 1pt
\usepackage{xcolor}
\usepackage{fvextra}
\usepackage{tcolorbox}
\usepackage{hyperref}

\definecolor{codebg}{HTML}{F4F6F8}
\definecolor{accent}{HTML}{2B5F8A}
\fvset{numbers=left, numbersep=6pt, fontsize=\footnotesize, frame=single,
  rulecolor=\color{gray!50}, breaklines=true, breakanywhere=true, xleftmargin=18pt}
\setlength{\parindent}{0pt}
\setlength{\parskip}{4pt}
\newcommand{\code}[1]{\texttt{#1}}
\newcommand{\req}[1]{\begin{tcolorbox}[colback=accent!6,colframe=accent!60,boxrule=0.5pt,left=4pt,right=4pt,top=2pt,bottom=2pt]\textbf{What it does:} #1\end{tcolorbox}}
\newcommand{\snap}[3]{\par\medskip\textbf{\code{#1}}, lines #2 --- #3}
% \promptbox{date}{prompt (Thai)}{English translation}{what the AI did}
\newcommand{\promptbox}[4]{%
\begin{tcolorbox}[colback=codebg,colframe=gray!50,boxrule=0.4pt,left=4pt,right=4pt,top=2pt,bottom=2pt]
{\small #1 \quad Claude Code (Claude Opus 5.5)}\par
\textbf{Prompt:} #2\par
\textit{English:} #3\par
\textbf{AI output:} #4
\end{tcolorbox}}

\title{\textbf{Part 3 -- Special Features: Code Snapshots and Generative AI Prompts}\\[2pt]
\large Simple Chat Application (Socket Programming)}
\author{{[Team name]} \\ {[Member names]}}
\date{}

\begin{document}
\maketitle

\section*{How to read this document}
We built four special features: an \textbf{AI word-guessing game} (our AI feature), a
\textbf{rock-paper-scissors game}, \textbf{sending files and pictures}, and a \textbf{Discord bridge}.
For each feature we give what it does, how it works over the network, \textbf{code snapshots} copied
from the project (with the real line numbers), and the \textbf{Generative AI prompts} we used, copied
word for word from our prompt log (\code{PromptRecord.md}) in the original Thai with an English translation.

\subsection*{How the features plug into the core}
All feature code is kept apart from the core chat (R1--R11). The server core calls the features only
through five hook functions, so a feature can be added or removed without changing the core. Every
feature still uses the same Socket.IO connection, and chat-related features reuse the core's permission
rules (R7, R11).

\snap{server/features/index.js}{29--59}{Feature registry and the five hooks}
\begin{Verbatim}[firstnumber=29]
const features = [cloud, discord, ping, game, bot, file];

// เรียกจาก core ครั้งเดียว หลังสร้าง io และ state (clients, groups)
export function initFeatures(app) {
  for (const f of features) f.init?.(app);
}

// เรียกจาก HTTP handler ของ core: feature ตัวแรกที่ return true เป็นผู้ตอบ request นั้น
// ไม่มีใครรับ → core ตอบ 404 เหมือนเดิม (R2: ไม่มี REST สำหรับข้อความแชท)
export function handleHttp(req, res) {
  return features.some((f) => f.handleHttp?.(req, res) === true);
}

// เรียกจาก io.on("connection") ของ core ทุกครั้งที่มี socket ใหม่
export function setupSocket(socket, ctx) {
  for (const f of features) f.setupSocket?.(socket, ctx);
}

// เรียกจาก core หลังส่ง group_msg ให้สมาชิกแล้ว (msg = { group, from, text, timestamp })
export function onGroupMessage(msg) {
  for (const f of features) f.onGroupMessage?.(msg);
}

// เรียกจาก register ของ core ก่อนจองชื่อ: feature ตัวแรกที่ปฏิเสธชนะ (เช่น bot สงวนชื่อ "AI")
export function checkRegister(name, payload) {
  for (const f of features) {
    const denied = f.checkRegister?.(name, payload);
    if (denied) return denied;
  }
  return null;
}
\end{Verbatim}

We asked for this structure before building any feature:

\promptbox{2026-10-28}{ก่อนทำเราแยกส่วนที่เป็น Feature เสริมได้ไหม อย่าง Server/Client ก็มี folder เพิ่ม handbook.md ก็มีเพิ่ม แล้วก็มีการเก็บ Test เฉพาะ Feature เสริมด้วย}{Before we start, can we separate the special features? For example, extra folders in Server/Client, an extra handbook, and separate tests only for the special features.}{Created the feature registry \code{server/features/index.js} with hooks, \code{client/src/features/}, \code{server/tests/features/} and \code{HandbookFeatures.md}.}

Every change was checked by our team: we ran the automated tests (\code{cd server \&\& npm test}),
tested manually on separate computers, and saved screenshots as evidence.

\clearpage
\section{Feature 1 -- AI word-guessing game}
\req{An AI chooses a secret English word (a person, an animal or an object) and posts hints
in a group, one by one. Group members race to guess it by typing in the group chat; the first correct guess
wins, and after 5 minutes nobody wins. A wrong guess gets a short AI hint comparing it with the answer.}

The AI is a \textbf{separate client program} (\code{ai-bot/}) that logs in to our server as
the reserved user \code{AI} with a secret token. A player's \code{/guess start} command goes to the server,
which checks group membership and forwards it to the bot as \code{bot\_request} (with an ack and a 10-second
timeout). The bot joins the group itself (R10) and sends every hint and announcement as a normal
\code{group\_msg} (R11), so all AI messages travel through the socket like a human's. The bot calls
Claude, ChatGPT or Gemini over HTTPS with a JSON schema (structured output). The API keys and the
answer stay on the bot only; the server and the clients never see them.

\subsection*{Code snapshots}
\snap{server/features/bot.js}{29--36}{Server: the name "AI" is reserved for the bot with the right token}
\begin{Verbatim}[firstnumber=29]
  // core เรียกตอน register ก่อนจองชื่อ (R3): ชื่อ "AI" (ตัวพิมพ์ใดก็ได้) ต้องมี token ตรงกับ AI_BOT_TOKEN
  // server ไม่ได้ตั้ง AI_BOT_TOKEN → ไม่มีใครใช้ชื่อ AI ได้เลย
  checkRegister(name, payload) {
    if (!isBotName(name)) return null;
    const token = process.env.AI_BOT_TOKEN;
    if (token && payload?.token === token) return null;
    return { code: "NAME_RESERVED", message: `Name "${name}" is reserved for the AI bot` };
  },
\end{Verbatim}

\snap{server/features/bot.js}{38--64}{Server: check membership, forward /guess to the bot, return its ack}
\begin{Verbatim}[firstnumber=38]
  setupSocket(socket, { requireName }) {
    for (const event of RELAYED) {
      socket.on(event, async (payload = {}, ack) => {
        const reply = (data) => typeof ack === "function" && ack(data);
        const me = requireName();
        if (!me) return;
        const group = payload?.group;
        // ต้องเป็นสมาชิกกลุ่ม (หลักเดียวกับ R11) คนนอกจึงสั่ง / ดูคำใบ้ของกลุ่มอื่นไม่ได้
        if (event !== "guess_ai") {
          const members = groups.get(group);
          if (event === "guess_start" && !members) return reply({ ok: false, code: "GROUP_NOT_FOUND" });
          if (event === "guess_start" && !members.has(me)) return reply({ ok: false, code: "NOT_A_MEMBER" });
          if (!members?.has(me)) return reply({ ok: false, code: "NO_GAME" });
        }

        const bot = clients.get(BOT_NAME);
        if (!bot) return reply({ ok: false, code: "BOT_OFFLINE", message: "the AI bot is offline (start it with: npm run bot)" });
        try {
          // from มาจาก server เสมอ (ไม่ใช้ค่าจาก payload) bot จึงเชื่อได้ว่าใครสั่ง
          const req = { event, from: me, group, category: payload?.category, ai: payload?.ai };
          reply(await bot.timeout(BOT_TIMEOUT_MS).emitWithAck("bot_request", req));
        } catch {
          reply({ ok: false, code: "BOT_TIMEOUT", message: "the AI bot did not answer" });
        }
      });
    }
  },
\end{Verbatim}

\snap{ai-bot/index.js}{26--34}{Bot: connect as a Socket.IO client; speak with group\_msg}
\begin{Verbatim}[firstnumber=26]
// reconnection: true (ต่างจาก client ของผู้ใช้) → server restart แล้ว bot กลับมาเอง
const socket = io(SERVER_URL, { transports: ["websocket"], reconnection: true });
let memberOf = new Set(); // กลุ่มที่ bot เป็นสมาชิกอยู่ (อ่านจาก group_list ที่ server broadcast, R9)

// ---- ส่งเข้า server ----
function say(group, text) {
  console.log(`[send] group_msg #${group}: ${text}`);
  socket.emit("group_msg", { group, text });
}
\end{Verbatim}

\snap{ai-bot/index.js}{61--65}{Bot: register as "AI" with the token}
\begin{Verbatim}[firstnumber=61]
// ---- รับจาก server ----
socket.on("connect", () => {
  console.log(`[bot] connected to ${SERVER_URL} — registering as "${BOT_NAME}"`);
  socket.emit("register", { name: BOT_NAME, token: TOKEN });
});
\end{Verbatim}

\snap{ai-bot/index.js}{82--98}{Bot: handle forwarded commands and receive guesses}
\begin{Verbatim}[firstnumber=82]
// คำสั่ง /guess ที่ server ส่งต่อมา: { event, from, group, category, ai } → ตอบผ่าน ack (server ส่งต่อให้คนสั่ง)
socket.on("bot_request", async (req, ack) => {
  console.log(`[recv] bot_request ${req.event} from ${req.from}${req.group ? ` #${req.group}` : ""}`);
  let reply;
  try {
    reply = await game.handleRequest(req);
  } catch (err) {
    reply = { ok: false, code: "BOT_ERROR", message: err.message };
  }
  if (typeof ack === "function") ack(reply);
});

// ข้อความในกลุ่มที่ bot เป็นสมาชิก = คำทาย (ข้ามข้อความของ bot เอง)
socket.on("group_msg", (msg) => {
  if (msg.from === BOT_NAME && !msg.source) return;
  game.onGroupMessage(msg);
});
\end{Verbatim}

\snap{ai-bot/guess-ai.js}{20--34}{JSON schema the AI must answer with}
\begin{Verbatim}[firstnumber=20]
// รูปแบบ JSON ที่บังคับให้ AI ตอบ (ใช้กับทุกเจ้า)
const PUZZLE_SCHEMA = {
  type: "object",
  properties: {
    word: { type: "string", description: "the secret answer in English, lowercase" },
    aliases: {
      type: "array",
      items: { type: "string" },
      description: "other accepted English answers: alternative English names, spellings, plural; lowercase; English only",
    },
    hints: { type: "array", items: { type: "string" }, description: `exactly ${HINT_COUNT} hints in English, from vague to very specific` },
  },
  required: ["word", "aliases", "hints"],
  additionalProperties: false,
};
\end{Verbatim}

\snap{ai-bot/ai-providers/claude.js}{15--39}{Calling Claude over HTTPS (one file per AI provider)}
\begin{Verbatim}[firstnumber=15]
  // คืน object ที่ parse จาก JSON ตาม schema แล้ว (structured output)
  async generateJson(prompt, schema) {
    client ??= new Anthropic();
    let response;
    try {
      response = await client.beta.messages.create({
        model: MODEL,
        max_tokens: 16000,
        betas: ["server-side-fallback-2026-07-01"],
        fallbacks: "default", // ถ้าโมเดลปฏิเสธ ให้ server ของ Anthropic ลองโมเดลสำรองให้
        output_config: { effort: "low", format: { type: "json_schema", schema } },
        messages: [{ role: "user", content: prompt }],
      });
    } catch (err) {
      if (err instanceof Anthropic.AuthenticationError) throw new Error("Claude API key is invalid (check ANTHROPIC_API_KEY)");
      if (err instanceof Anthropic.RateLimitError) throw new Error("Claude is busy right now, try again in a minute");
      if (err instanceof Anthropic.APIConnectionError) throw new Error("cannot reach Claude (check internet connection)");
      if (err instanceof Anthropic.APIError) throw new Error(`Claude error ${err.status}`);
      throw err;
    }
    if (response.stop_reason === "refusal") throw new Error("Claude declined to make a puzzle, try another category");
    const text = response.content.find((b) => b.type === "text")?.text;
    if (!text) throw new Error("Claude returned no puzzle");
    return JSON.parse(text);
  },
\end{Verbatim}

\snap{ai-bot/game.js}{178--199}{Check each guess; wrong guess -> AI feedback in the group}
\begin{Verbatim}[firstnumber=178]
// ทุกข้อความแชทในกลุ่มที่ bot เป็นสมาชิก (index.js ตัดข้อความของ bot เองออกแล้ว) = คำทาย ถ้ากลุ่มนั้นมีเกมอยู่
export function onGroupMessage(msg) {
  const g = games.get(msg.group);
  if (!g || g.status !== "running") return;
  if (g.answers.has(normalizeGuess(msg.text))) {
    announce(msg.group, g.host, `🎉 ${msg.from} got it in ${elapsed(g)}! The answer was ${g.word.toUpperCase()}.`);
    return end(msg.group);
  }

  // ทายผิด → ให้ AI ใบ้เทียบคำทายกับคำตอบจริง (ไม่เฉลย) แล้วประกาศในกลุ่ม
  if (!looksLikeGuess(msg.text) || g.feedbackCount >= MAX_FEEDBACK_PER_GAME) return;
  const now = Date.now();
  if (now - (g.lastFeedback.get(msg.from) ?? 0) < FEEDBACK_COOLDOWN_MS) return;
  g.lastFeedback.set(msg.from, now);
  g.feedbackCount++;

  const guess = msg.text.trim();
  generateFeedback(g, guess, g.aiProvider, g.hints.slice(0, g.shown)).then((feedback) => {
    if (games.get(msg.group) !== g) return; // มีคนถูก / หมดเวลา / หยุดไปแล้วระหว่างรอ AI → ทิ้ง
    announce(msg.group, g.host, `💬 ${msg.from} guessed "${guess}" — ${feedback ?? "not it, keep trying!"}`);
  });
}
\end{Verbatim}

\subsection*{Generative AI prompts used}
\promptbox{2026-10-28}{อีก Feature ทำเป็นเกมให้ AI คิดคำมาหนึ่งคำ หลังจากนั้นให้ AI ใบ้ คน สัตว์ สิ่งของ โดยการ ส่งข้อความ Socket มาเรื่อยๆ 5 นาที ถือว่าไม่มีใครชนะใครทายถูกก่อนหหน้านั้นคนนั้นชนะ แล้วให้ คนในกลุ่มแข่งกันทาย}{Another feature: a game where the AI thinks of one word, then gives hints (person, animal, object) by sending socket messages over and over. After 5 minutes nobody wins; whoever guesses correctly before that wins. Members of the group compete to guess.}{Built the game with the Claude API (structured JSON output: word, aliases, 10 hints), one game per group, hints spread over 5 minutes, slash commands \code{/guess start|hints|stop}, an \code{[ai]} tag in the UI, and 24 automated tests with a fake AI.}
\promptbox{2026-10-28}{ลืมบอกไปว่าทายโดยใช้อค่ English กับ เพิ่ม API ของ AI ตัวอื่นไปด้วยอย่าง chatGPT gemini}{I forgot to say that guesses are in English only. Also add other AI APIs such as ChatGPT and Gemini.}{Split the AI code into one file per provider (Claude, OpenAI, Gemini), all with \code{generateJson(prompt, schema)}; answers and aliases are English only; \code{/guess start [category] [ai]} chooses the AI.}
\promptbox{2026-10-28}{อยากให้เวลาทุกครั้งที่ทายผิดให้ AI ช่วยใบ้อ้างอิงจากคำตอบของ client เทียบกับคำตอบจริงแล้วพิมพ์ลงกลุ่ม}{Every time someone guesses wrong, I want the AI to give a hint comparing the client's guess with the real answer, and post it in the group.}{Added \code{generateFeedback} (the AI must not reveal the answer; replies that leak it are dropped), with limits: English guesses of 1--4 words, a 5-second cooldown per player and at most 40 AI replies per game.}
\promptbox{2026-10-28}{อยากให้เวลา AI ตอบให้ส่ง Socket ให้ Server ด้วย \newline (เลือก: AI เป็น client อีกตัว)}{When the AI answers, I want it to send through the socket to the server too. (Chosen: the AI is another client.)}{Moved the AI into a separate program \code{ai-bot/} that connects with \code{socket.io-client} as user \code{AI}; the server only reserves the name and relays commands (\code{server/features/bot.js}) and no longer holds any API key.}

\clearpage
\section{Feature 2 -- Rock-paper-scissors game}
\req{Two online users play rock-paper-scissors (best of 1, 3 or 5) with slash commands
(\code{/play rps <name>}, \code{/accept}, \code{/move r|p|s}, \code{/quit}). Game messages appear in their
private room.}

We designed a small \textbf{application-level protocol} on Socket.IO. Every command is an
event with an ack \code{\{ok, code\}}. The server is authoritative: it keeps the state
(\code{pending} $\rightarrow$ \code{playing} $\rightarrow$ \code{finished}), checks the rules and decides the
winner. The first move of a round is kept secret at the server; the opponent only learns that a move was
made. When both have chosen, the server reveals both moves at once. Updates go only to the two players
(like R7), and a disconnect ends the game.

\subsection*{Code snapshots}
\snap{server/features/game.js}{1--10}{Message flow of the game protocol}
\begin{Verbatim}[firstnumber=1]
// Feature: game — เกมผลัดกันเล่นผ่าน socket (เริ่มที่ เป่ายิ้งฉุบ: Rock-Paper-Scissors)
//
//   C1 ──game_challenge──> server ──game_update (challenge)──> C1, C2
//   C2 ──game_respond────> server ──game_update (accepted)───> C1, C2
//   C1 ──game_move rock──> server  (เก็บไว้เป็นความลับ)  ──game_update (moved: แค่บอกว่า C1 เลือกแล้ว)──> C1, C2
//   C2 ──game_move paper─> server  ตัดสิน ──game_update (round/finished: เปิดคำตอบทั้งคู่ + คะแนน)──> C1, C2
//
// Server เป็นคนตัดสินทั้งหมด (server-authoritative): client แค่ส่ง "การกระทำ" server ตรวจว่าถูกกติกาไหม
// และคำตอบของคนที่เลือกก่อน**ไม่ถูกส่งให้ใคร**จนกว่าจะครบ 2 คน ถ้าส่งตรงหากัน อีกฝ่ายจะเห็นก่อนแล้วโกงได้
// update ทุกอันส่งให้ผู้เล่น 2 คนเท่านั้น (หลักเดียวกับ R7)
\end{Verbatim}

\snap{server/features/game.js}{35--53}{Public state (no hidden moves) is sent to the 2 players only}
\begin{Verbatim}[firstnumber=35]
// สถานะที่ส่งให้ผู้เล่น: **ไม่มี** คำตอบที่ยังไม่เปิด มีแค่ชื่อคนที่เลือกแล้ว
function publicState(g) {
  return {
    id: g.id,
    game: g.game,
    title: GAMES[g.game],
    players: g.players,
    bestOf: g.bestOf,
    status: g.status,
    round: g.round,
    score: { ...g.score },
    moved: g.players.filter((p) => p in g.moves),
  };
}

function send(g, event, extra = {}) {
  const update = { ...publicState(g), event, ...extra };
  for (const p of g.players) clients.get(p)?.emit("game_update", update);
}
\end{Verbatim}

\snap{server/features/game.js}{129--161}{game\_move: keep the move secret until both chose, then decide}
\begin{Verbatim}[firstnumber=129]
    socket.on("game_move", ({ id, move } = {}, ack) => {
      const me = requireName();
      if (!me) return;
      const g = myGame(id, me);
      if (!g) return reply(ack, { ok: false, code: "NO_GAME" });
      if (g.status !== "playing") return reply(ack, { ok: false, code: "NOT_PLAYING" });
      const m = normalizeMove(move);
      if (!m) return reply(ack, { ok: false, code: "BAD_MOVE" });
      if (me in g.moves) return reply(ack, { ok: false, code: "ALREADY_MOVED" });

      g.moves[me] = m;
      reply(ack, { ok: true, move: m }); // บอกคนเลือกคนเดียวว่าเลือกอะไร

      if (Object.keys(g.moves).length < 2) return send(g, "moved", { by: me }); // อีกฝ่ายรู้แค่ว่า "เลือกแล้ว"

      // ครบ 2 คน → เปิดคำตอบพร้อมกัน + ตัดสิน
      const [a, b] = g.players;
      const w = rpsWinner(g.moves[a], g.moves[b]);
      const result = { round: g.round, moves: { ...g.moves }, winner: w === 0 ? null : g.players[w - 1] };
      if (result.winner) g.score[result.winner]++;
      g.moves = {};
      g.round++;

      const need = Math.ceil(g.bestOf / 2);
      const champion = g.players.find((p) => g.score[p] >= need);
      if (champion) {
        g.status = "finished";
        send(g, "finished", { result, winner: champion });
        end(g);
      } else {
        send(g, "round", { result });
      }
    });
\end{Verbatim}

\snap{server/features/game.js}{174--182}{A player disconnects -> the game is abandoned}
\begin{Verbatim}[firstnumber=174]
    // ผู้เล่นหลุดการเชื่อมต่อ → จบเกม แจ้งอีกฝ่าย (core ลบชื่อออกจาก clients ไปก่อนแล้ว จึงส่งถึงคนที่เหลือคนเดียว)
    socket.on("disconnect", () => {
      const name = socket.data.name;
      const g = name && games.get(playing.get(name));
      if (!g) return;
      g.status = "ended";
      send(g, "abandoned", { by: name });
      end(g);
    });
\end{Verbatim}

\snap{client/src/features/game/index.js}{9--16}{Client: show game updates in the private room}
\begin{Verbatim}[firstnumber=9]
export function setupGame(socket, notify, getMe) {
  current = null; // login ใหม่ = เริ่มใหม่
  socket.on("game_update", (u) => {
    const me = getMe();
    current = isGameOver(u.event) ? null : u;
    for (const line of formatUpdate(u, me)) notify(opponentOf(u, me), line);
  });
}
\end{Verbatim}

\snap{client/src/features/game/index.js}{32--37}{Client: send an action and wait for the ack}
\begin{Verbatim}[firstnumber=32]
// ส่ง action แบบมี ack แล้วพิมพ์ error ถ้าไม่ผ่าน (คืน reply ถ้าผ่าน)
async function act(ctx, event, data, subject) {
  const reply = await ctx.socket.timeout(3000).emitWithAck(event, data);
  if (!reply.ok) ctx.print(`error: ${ERRORS[reply.code]?.(subject) ?? reply.code}`);
  return reply.ok ? reply : null;
}
\end{Verbatim}

\subsection*{Generative AI prompts used}
\promptbox{2026-10-28}{สามารถใช้ Socket ทำเกมที่ใช้ Text ตอบโต้กันสองฝั่งได้ไหม \newline เริ่มที่ RPS ก่อนเลย}{Can we use the socket to make a text game where both sides respond to each other? Start with RPS.}{Built \code{server/features/game.js} (server-authoritative state and rules, hidden moves, updates only to the 2 players, best of 1/3/5, decline / quit / 60-second expiry / abandoned on disconnect) and the client commands; 31 automated tests.}

\clearpage
\section{Feature 3 -- Sending files and pictures}
\req{Users send any file (up to 5\,MB) to a private chat or a group with a \code{[+file]}
button, drag-and-drop or paste. Pictures are previewed in the chat, and every file has a progress bar and a
download link.}

The file is split into \textbf{64\,KB binary chunks}. Each chunk is one Socket.IO event with
framing \code{\{id, seq, data\}}; the data is sent as binary, not base64. The sender waits for the ack of a
chunk before sending the next one (\textbf{stop-and-wait}, flow control at the application level, similar to
TCP ACKs). The server checks the receiver with the same rules as chat messages (R7: online user, R11: group
member), checks the owner, order and size of each chunk, and forwards it only to the receivers. It never
stores the file. The receiver joins the chunks into a file at the end.

\subsection*{Code snapshots}
\snap{server/features/file.js}{12--16}{Limits: chunk size, file size, idle timeout}
\begin{Verbatim}[firstnumber=12]
const CHUNK_SIZE = 64 * 1024; // 64 KB ต่อชิ้น (ต่ำกว่า maxHttpBufferSize 1 MB ของ Socket.IO มาก)
const MAX_FILE_SIZE = Number(process.env.FILE_MAX_BYTES) || 5 * 1024 * 1024; // 5 MB
const IDLE_TIMEOUT_MS = Number(process.env.FILE_IDLE_TIMEOUT_MS) || 30_000; // ผู้ส่งเงียบนานเกิน → ยกเลิก
const MAX_ACTIVE_PER_USER = 3;
const MAX_NAME_LENGTH = 100;
\end{Verbatim}

\snap{server/features/file.js}{32--44}{Forward only to the receiver or the group members}
\begin{Verbatim}[firstnumber=32]
// ส่งถึง "ผู้รับ" (ไม่รวมผู้ส่ง): private = อีกฝ่าย 1 คน / group = สมาชิกทุกคนใน room ยกเว้นผู้ส่ง
function toReceivers(t, event, data) {
  if (t.group) t.socket.to(`group:${t.group}`).emit(event, data);
  else t.target.emit(event, data);
}
// ส่งถึงทุกคนที่เกี่ยวข้อง (ผู้รับ + ผู้ส่ง)
function toAll(t, event, data) {
  if (t.group) io.to(`group:${t.group}`).emit(event, data);
  else {
    t.target.emit(event, data);
    t.socket.emit(event, data);
  }
}
\end{Verbatim}

\snap{server/features/file.js}{90--103}{file\_offer: the same permission checks as R7 / R11}
\begin{Verbatim}[firstnumber=90]
      // ตรวจผู้รับด้วยกฎเดียวกับข้อความแชท: private (R7) ต้องออนไลน์ / group (R11) ต้องเป็นสมาชิก
      const t = { socket, from: me, name, size, type: String(type || "application/octet-stream").slice(0, 100) };
      if (group !== undefined) {
        const members = groups.get(group);
        if (!members) return reply(ack, { ok: false, code: "GROUP_NOT_FOUND", message: "Group not found" });
        if (!members.has(me)) return reply(ack, { ok: false, code: "NOT_A_MEMBER", message: `You are not a member of "${group}"` });
        t.group = group;
      } else {
        const target = clients.get(to);
        if (!target) return reply(ack, { ok: false, code: "USER_NOT_FOUND", message: `User "${to}" is not online` });
        if (target === socket) return reply(ack, { ok: false, code: "SELF", message: "Cannot send a file to yourself" });
        t.to = to;
        t.target = target;
      }
\end{Verbatim}

\snap{server/features/file.js}{119--133}{file\_chunk: check owner, order and size, then forward}
\begin{Verbatim}[firstnumber=119]
    socket.on("file_chunk", ({ id, seq, data } = {}, ack) => {
      if (!requireName()) return;
      const t = own(id);
      if (!t) return reply(ack, { ok: false, code: "NO_TRANSFER" });
      // ชิ้นต้องมาตามลำดับ 0, 1, 2, ... (ผู้รับเอาไปต่อกันตามลำดับที่ได้รับ)
      if (seq !== t.nextSeq) return reply(ack, { ok: false, code: "BAD_SEQ", expected: t.nextSeq });
      if (!Buffer.isBuffer(data) || data.length === 0 || data.length > CHUNK_SIZE || t.received + data.length > t.size) {
        return reply(ack, { ok: false, code: "BAD_CHUNK" });
      }
      t.nextSeq++;
      t.received += data.length;
      touch(t);
      toReceivers(t, "file_chunk", { id: t.id, seq, data });
      reply(ack, { ok: true, received: t.received });
    });
\end{Verbatim}

\snap{client/src/features/file/transfer.js}{64--95}{Client: send chunk by chunk, waiting for each ack}
\begin{Verbatim}[firstnumber=64]
// target = { to: "C2" } หรือ { group: "g1" } / print = แสดงข้อความ error เฉพาะจอเรา
export async function sendFile(socket, target, file, print) {
  const ref = Math.random().toString(36).slice(2);
  const out = { url: URL.createObjectURL(file), startedAt: performance.now(), cancelled: false };
  outgoing.set(ref, out);
  let id = null;
  try {
    const offer = await socket
      .timeout(5000)
      .emitWithAck("file_offer", { ...target, name: file.name || "file", size: file.size, type: file.type, ref });
    if (!offer.ok) return print(`error: ${offer.message ?? offer.code}`);
    id = offer.id;

    // ส่งทีละชิ้น: อ่านเฉพาะส่วนของไฟล์ (slice) → ArrayBuffer → emit แบบ binary → รอ ack ก่อนชิ้นถัดไป
    for (let seq = 0, pos = 0; pos < file.size; seq++, pos += offer.chunkSize) {
      if (out.cancelled) return; // server ยกเลิก (เช่น ผู้รับออฟไลน์)
      const data = await file.slice(pos, pos + offer.chunkSize).arrayBuffer();
      const r = await socket.timeout(10_000).emitWithAck("file_chunk", { id, seq, data });
      if (!r.ok) throw new Error(r.code);
      const t = transfers.get(id);
      if (t) ui.update(t.room, id, { done: r.received });
    }
    if (out.cancelled) return;
    const end = await socket.timeout(5000).emitWithAck("file_end", { id });
    if (!end.ok) throw new Error(end.code);
  } catch (err) {
    if (id && !out.cancelled) socket.emit("file_cancel", { id }); // บอก server + ผู้รับให้ทิ้งไฟล์นี้
    print(`error: could not send ${file.name} (${err.message === "operation has timed out" ? "timed out" : err.message})`);
  } finally {
    outgoing.delete(ref);
  }
}
\end{Verbatim}

\snap{client/src/features/file/transfer.js}{33--53}{Client: collect chunks and build the file}
\begin{Verbatim}[firstnumber=33]
  // ชิ้นไฟล์มาตามลำดับ (server ตรวจ seq แล้ว) เก็บต่อท้ายไว้ก่อน
  socket.on("file_chunk", ({ id, data }) => {
    const t = transfers.get(id);
    if (!t || t.out) return;
    t.chunks.push(data);
    t.received += data.byteLength;
    ui.update(t.room, id, { done: t.received });
  });

  socket.on("file_end", ({ id }) => {
    const t = transfers.get(id);
    if (!t) return;
    transfers.delete(id);
    if (t.out) {
      const sec = (performance.now() - t.out.startedAt) / 1000;
      return ui.update(t.room, id, { status: "done", done: t.size, seconds: sec });
    }
    // ประกอบทุกชิ้นเป็นไฟล์ / ชนิดไฟล์มาจากผู้ส่ง (เชื่อไม่ได้) จึงใช้ตามจริงเฉพาะรูปที่แสดงได้ปลอดภัย
    const blob = new Blob(t.chunks, { type: canPreview(t.type) ? t.type : "application/octet-stream" });
    ui.update(t.room, id, { status: "done", done: t.size, url: URL.createObjectURL(blob) });
  });
\end{Verbatim}

\subsection*{Generative AI prompts used}
\promptbox{2026-10-28}{คิดว่า ping ทำอะไรได้น้อยไปไหม \newline เปลี่ยนเป็น ส่งไฟล์ส่งรูป}{Do you think ping does too little? Change it to sending files and pictures.}{Built \code{server/features/file.js} (offer / chunk / end / cancel, permission checks like R7 / R11, size limits, 30-second idle timeout, no storage) and the client (chunked sending with stop-and-wait, progress bar, image preview, download); 27 automated tests. Tested with a 300\,KB file between 2 processes: the SHA-256 hashes matched.}

\clearpage
\section{Feature 4 -- Discord bridge}
\req{The group \code{\#discord} in our app is connected to a channel in a real Discord server.
Messages typed in either place appear in the other.}

Our server acts as a \textbf{client of Discord}. It logs in through the Discord Gateway
(WebSocket over TLS) to receive new messages, and posts messages with the Discord REST API (HTTPS).
Outgoing: after the core has checked membership and delivered a \code{group\_msg} (R11), the
\code{onGroupMessage} hook posts it to Discord. Incoming: a Discord message is converted to a
\code{group\_msg} with \code{source: "discord"} and emitted only to the group's room. Messages from bots are
ignored, so messages never loop. The token is read from \code{server/.env}; without it the feature turns
itself off and the app still works.

\subsection*{Code snapshots}
\snap{server/features/discord.js}{24--46}{Convert messages in both directions; ignore bots}
\begin{Verbatim}[firstnumber=24]
// ---- แปลงข้อความ (pure function ใช้ใน test ได้โดยไม่ต้องต่อ Discord) ----

// ขาออก: ข้อความกลุ่มในแอป → ข้อความใน Discord เช่น "**C1**: hi"
export function toDiscordText(msg) {
  return `**${escapeMarkdown(msg.from)}**: ${msg.text}`;
}

// ขาเข้า: Discord message → payload group_msg ของแอป (source: "discord" ให้ client แสดงป้ายกำกับ)
// คืน null ถ้าไม่ควรส่งต่อ (ข้อความจาก bot รวมถึงตัวเอง กันวนลูป / ข้อความว่าง)
export function fromDiscordMessage(message, group = GROUP) {
  if (message.author?.bot) return null;
  const parts = [message.cleanContent, ...[...(message.attachments?.values?.() ?? [])].map((a) => a.url)];
  let text = parts.filter(Boolean).join("\n").trim();
  if (!text) return null;
  if (text.length > MAX_TEXT_LENGTH) text = text.slice(0, MAX_TEXT_LENGTH - 1) + "…";
  return {
    group,
    from: message.member?.displayName ?? message.author?.globalName ?? message.author?.username ?? "unknown",
    text,
    timestamp: (message.createdAt ?? new Date()).toISOString(),
    source: "discord",
  };
}
\end{Verbatim}

\snap{server/features/discord.js}{55--65}{Create the \#discord group and connect to the Gateway}
\begin{Verbatim}[firstnumber=55]
  init({ io, groups }) {
    if (!enabled) return;

    // กลุ่มที่ bridge สร้างไว้ตั้งแต่เริ่ม แบบไม่มีสมาชิก ใครอยากคุยก็ [join] เอง (R10)
    // (feature แตะ groups ได้ที่นี่ที่เดียว ตอนเริ่ม server ยังไม่มีใครสร้างกลุ่มชื่อนี้)
    if (!groups.has(GROUP)) groups.set(GROUP, new Set());

    console.log(`[discord] connecting to Discord Gateway... (group #${GROUP}, channel ${CHANNEL_ID})`);
    const client = new Client({
      intents: [GatewayIntentBits.Guilds, GatewayIntentBits.GuildMessages, GatewayIntentBits.MessageContent],
    });
\end{Verbatim}

\snap{server/features/discord.js}{76--87}{Incoming: Discord -> members of the group only}
\begin{Verbatim}[firstnumber=76]
    // ขาเข้า: ข้อความใหม่ใน channel → ส่งให้สมาชิกกลุ่มในแอป (เฉพาะ room ของกลุ่ม เหมือน R11)
    client.on(Events.MessageCreate, (message) => {
      if (message.channelId !== CHANNEL_ID || message.author?.bot) return;
      const msg = fromDiscordMessage(message);
      if (!msg) {
        // มีคนพิมพ์แต่เนื้อหาว่าง: เกือบทุกครั้งเกิดจากยังไม่ได้เปิด MESSAGE CONTENT INTENT ใน Developer Portal
        console.warn(`[discord] got an empty message from ${message.author?.username} — enable MESSAGE CONTENT INTENT (Developer Portal → Bot)`);
        return;
      }
      console.log(`[discord] in  <- ${msg.from}: ${msg.text}`);
      io.to(`group:${GROUP}`).emit("group_msg", msg);
    });
\end{Verbatim}

\snap{server/features/discord.js}{120--133}{Outgoing: group message -> Discord REST API}
\begin{Verbatim}[firstnumber=120]
  // ขาออก: ข้อความกลุ่ม (ผ่านการตรวจสิทธิ์ของ core แล้ว) → โพสต์ใน Discord
  // ข้อความที่มาจาก Discord ไม่ผ่าน core จึงไม่ถูกส่งกลับไปซ้ำ (ไม่วนลูป)
  onGroupMessage(msg) {
    if (!enabled || msg.group !== GROUP) return;
    if (!channel) {
      // bot ยังต่อ Discord ไม่สำเร็จ ไม่ข้ามเงียบ ๆ ให้ดูสาเหตุจาก log ตอนเริ่ม (login failed / cannot open channel)
      console.warn(`[discord] not connected to Discord — message from ${msg.from} NOT forwarded (see [discord] startup logs)`);
      return;
    }
    console.log(`[discord] out -> ${msg.from}: ${msg.text}`);
    channel
      .send({ content: toDiscordText(msg), allowedMentions: { parse: [] } }) // ไม่ ping @everyone / @user
      .catch((err) => console.error(`[discord] send failed: ${err.message}`));
  },
\end{Verbatim}

\snap{server/index.js}{187--189}{Core hook: after the group message is delivered}
\begin{Verbatim}[firstnumber=187]
    const msg = { group, from, text, timestamp: now() };
    io.to(`group:${group}`).emit("group_msg", msg);
    onGroupMessage(msg); // จุดเชื่อม Feature เสริม: เช่น discord bridge ส่งต่อข้อความออกไป Discord
\end{Verbatim}

\subsection*{Generative AI prompts used}
\promptbox{2026-10-28}{อยากทำให้เกี่ยวข้องกับ Discord \newline (เลือก: เชื่อมกับ Discord จริง --- bridge)}{I want to make something related to Discord. (Chosen: connect to real Discord --- a bridge.)}{Built \code{server/features/discord.js} with discord.js (Gateway in, REST out, loop prevention, turns itself off without a token), the \code{onGroupMessage} hook in the core, the \code{[discord]} tag in the client, a setup guide and 14 automated tests.}
\promptbox{2026-10-28}{ขั้นตอน 1.4 หาไม่เจอ}{I can't find step 1.4 (the channel ID).}{Explained two ways to get the channel ID (from the browser URL, or with Developer Mode) and updated the guide.}
\promptbox{2026-10-28}{พิมพ์หา Discord แล้วไม่ขึ้น Discord พิมพ์ก็ไม่ขึ้น Log ใน Render ล่าสุดคือการ Join ห้อง}{Messages typed to Discord don't show up, and messages typed in Discord don't show up either. The latest log on Render is the room join.}{Found the likely cause (the environment variables were not set on the host) and added clearer logs: connecting, \code{NOT forwarded} when the bot is not connected, and a warning to enable MESSAGE CONTENT INTENT.}
\promptbox{2026-10-28}{(วาง log จาก Render: \code{connecting to Discord Gateway...} แล้วไม่มี \code{logged in} / \code{login failed} ตามมา และส่งข้อความแล้วได้ \code{NOT forwarded})}{(Pasted the host log: \code{connecting to Discord Gateway...} with no \code{logged in} or \code{login failed} after it, and \code{NOT forwarded} when sending.)}{Diagnosed a stuck login (likely Discord rate-limiting the host's shared IP) and added step-by-step debug logs that never print the token. Running the server on our own machine worked, which confirmed the problem was the host.}

\end{document}
